% ============================================================================
% GEM SECTION 3 --- Experiments.
%
% REWRITTEN 2026-08-27 FROM TEXT SUPPLIED VERBATIM.  The prose below is
% transcribed as given; no sentence has been reworded, shortened or resequenced
% to fit the page limit.  Page fitting is the author's call, not this file's.
%
% The organising change from the previous version is that 3.2 now states the
% GuacaMol -> frozen R_theta -> ZINC transfer explicitly, as a property of the
% reference process rather than as an OOD-generalisation claim.
%
% TWO PLACEHOLDERS REMAIN, exactly as supplied.  They are NOT results:
%   [Insert matched 800-source result here.]
%   [Insert final frozen 30-cell x 3-run results here.]
% Table 1 carries the matching bracketed cells.  Nothing here may be presented
% as measured until those are filled from a checked artifact.
% ============================================================================

\section{Experiments}
\label{sec:gem-experiments}

Our experiments address three questions. First, we ask whether the frozen reference
law learns state-dependent preferences among executable molecular edits. Second,
we test whether a reference process learned once on GuacaMol transfers unchanged
to externally defined molecular-design benchmarks. Finally, we study
trajectory-level capabilities that endpoint benchmarks do not capture: using
downstream reachability to guide molecular edits, retargeting realized molecular
states when objectives change, and enforcing constraints throughout a
trajectory.

\subsection{The reference process learns state-dependent edit preferences}
\label{sec:gem-exp-rtheta}

The rewrite system specifies which molecular edits are executable from a given
state, but it does not determine how probability should be distributed among
them. We therefore test whether the frozen reference law $R_\theta$ learns
state-dependent transition preferences beyond the global frequencies of
different edit families.

We evaluate 3{,}545 held-out transitions whose source molecules are disjoint
from reference-law training. As a matched baseline, we assign probability
according to empirical edit-family frequencies from the training set while
preserving the same executable support and canonical successor set as
$R_\theta$. The comparison therefore isolates how probability is allocated among
the molecular transitions available from each state.

$R_\theta$ reduces canonical-successor NLL from 5.37 to 3.90 nats relative to
this empirical-family baseline. A matched decomposition shows that 86.7\% of the
improvement comes from learning which successor is appropriate for the current
molecule, rather than simply learning which edit families are common overall
(\cref{tab:rtheta-cells}). The learned signal is therefore primarily state
specific: the executor supplies the available molecular transitions, while
$R_\theta$ learns how preference should be distributed across them.

We keep $R_\theta$ fixed in all subsequent experiments and use it as a
goal-independent reference process for molecular design. Its transition
probabilities are not intended to represent physical kinetics or synthetic-route
likelihoods. Results restricted to edit families with data-backed transition
supervision are reported in \cref{app:exp-rtheta}.

\subsection{A molecular process learned once transfers to external design tasks}
\label{sec:gem-exp-external}

We next evaluate whether the molecular process learned on GuacaMol can be reused
without retraining $R_\theta$. We keep the reference process and executable
rewrite system fixed and apply COMPOSE to two external benchmarks with different
design objectives: similarity-constrained QED editing on ZINC-250k, where
molecules can grow or shrink during editing, and protein-specific lead
optimization, where molecular properties define feasibility and docking score
provides the optimization objective. Only the task-specific controller changes
between these settings.

\paragraph{Similarity-constrained QED editing.}
We first consider the ZINC-250k editing benchmark used by GrIDDD
\citep{ninniri2025griddd}. Source molecules have QED in $[0.7,0.8]$, and a source
is counted as solved if at least one of $K$ returned molecules reaches QED
$\geq 0.9$ while maintaining Morgan-fingerprint Tanimoto similarity $\geq 0.4$ to
the source. GrIDDD is a particularly relevant comparator because its generative
process can insert and delete nodes, allowing molecular size to change during
editing. Its reported model is trained directly on ZINC-250k.

For COMPOSE, the GuacaMol-trained $R_\theta$ from \cref{sec:gem-exp-rtheta} is
left unchanged. A QED-specific controller is fit on development sources disjoint
from evaluation and acts on the same executable molecular process. Each of the
$K$ returned candidates corresponds to one controlled trajectory. We count a
trajectory as successful when it first reaches the target region, and do not
retrospectively add intermediate molecules to the candidate set.

On the official 800-source panel, COMPOSE solves 446 of the 800 sources
(55.8\%, 95\% CI $[52.3,59.2]$) with $K=8$, compared with the 45.1\% success
reported by GrIDDD with $K=20$ (\cref{tab:external}). COMPOSE therefore solves
more sources while returning fewer than half as many candidate trajectories per
source. Among solved sources with at least two distinct non-source candidates,
the mean within-source pairwise Tanimoto distance is 0.4999.

\paragraph{Similarity-constrained lead optimization.}
We next evaluate the same frozen reference process on the GenMol T4 benchmark
\citep{lee2025genmol}, which combines molecular feasibility with a
protein-specific optimization objective. The benchmark contains five protein
targets, three seed molecules per target, and similarity thresholds
$\delta\in\{0.4,0.6\}$. Returned molecules must satisfy QED $\geq 0.6$, SA
$\leq 4$, and Morgan-fingerprint Tanimoto similarity $\geq\delta$ to the seed.
Among molecules satisfying these constraints, performance is measured by
QuickVina2 docking score, with lower scores indicating stronger predicted
binding.

We retain the same GuacaMol-trained $R_\theta$ and executable rewrite system used
above, replacing only the task-specific controller. Here QED, SA, and similarity
determine whether a candidate is eligible, while docking provides the
optimization signal. Because these feasibility checks are inexpensive, they are
evaluated before the protein oracle and docking calls are reserved for eligible
molecules.

\begin{table}[!ht]
\centering\footnotesize\setlength{\tabcolsep}{4pt}
{\renewcommand{\arraystretch}{1.25}%
\begin{tabular}{ll >{\columncolor{composeteal!15}}r rrr >{\columncolor{composeteal!15}}r rrr}
\toprule
& Seed & \multicolumn{4}{c}{$\delta=0.4$} & \multicolumn{4}{c}{$\delta=0.6$} \\
\cmidrule(lr){3-6}\cmidrule(lr){7-10}
Target & score & \textbf{COMPOSE} & GenMol & RetMol & GraphGA & \textbf{COMPOSE} & GenMol & RetMol & GraphGA \\
\midrule
PARP1 & -7.3 & \textbf{-10.7} & -10.6 & -9.0 & -8.3 & \textbf{-11.1} & -10.4 & --- & -8.6 \\
 & -7.8 & \textbf{-11.5} & -11.0 & -10.7 & -8.9 & \textbf{-10.9} & -9.7 & --- & -8.1 \\
 & -8.2 & \textbf{-11.6} & -11.3 & -10.9 & --- & \textbf{-10.2} & -9.2 & --- & --- \\
FA7 & -6.4 & \textbf{-10.2} & -8.4 & -8.0 & -7.8 & \textbf{-8.7} & -7.3 & -7.6 & -7.6 \\
 & -6.7 & \textbf{-8.9} & -8.4 & --- & -8.2 & -7.2 & \textbf{-7.6} & --- & \textbf{-7.6} \\
 & -8.5 & \textbf{-9.3} & --- & --- & --- & \textbf{-7.9} & --- & --- & --- \\
5HT1B & -4.5 & -11.7 & \textbf{-12.9} & -12.1 & -11.7 & \textbf{-13.4} & -12.1 & --- & -11.3 \\
 & -7.6 & -11.2 & \textbf{-12.3} & -9.0 & -12.1 & -11.1 & \textbf{-12.0} & -10.0 & \textbf{-12.0} \\
 & -9.8 & -11.4 & \textbf{-11.6} & --- & --- & --- & \textbf{-10.5} & --- & --- \\
BRAF & -9.3 & --- & \textbf{-10.8} & --- & -9.8 & --- & --- & --- & --- \\
 & -9.4 & -11.2 & -10.8 & \textbf{-11.6} & --- & \textbf{-10.6} & -9.7 & --- & --- \\
 & -9.8 & -10.9 & -10.6 & --- & \textbf{-11.6} & --- & \textbf{-10.5} & --- & -10.4 \\
JAK2 & -7.7 & \textbf{-10.3} & -10.2 & -8.2 & -8.7 & \textbf{-9.3} & \textbf{-9.3} & --- & -8.1 \\
 & -8.0 & -9.8 & \textbf{-10.0} & -9.0 & -9.2 & \textbf{-9.7} & -9.4 & --- & -9.2 \\
 & -8.6 & \textbf{-10.5} & -9.8 & --- & --- & \textbf{-10.2} & --- & --- & --- \\
\bottomrule
\end{tabular}}
\caption{\small\textbf{Similarity-constrained lead optimization (GenMol T4).} QuickVina2 docking score of the best feasible lead, lower is better, subject to QED $\geq0.6$, SA $\leq4$, and Tanimoto similarity $\geq\delta$ to the seed. Baseline values are from GenMol Table 4; bold marks the best entry in each cell. COMPOSE is feasible in 26/30 cells, matching GenMol, and improves the paired mean by $0.335$ kcal/mol over the 23 cells where both are feasible (95\% CI $0.019$--$0.651$), using 500 docking-oracle evaluations per cell against 3{,}000 for the published GenMol result.}
\label{tab:t4}
\end{table}


COMPOSE returns a feasible molecule in \TFourCoverage{} of the 30 benchmark
settings, matching GenMol's coverage (\cref{tab:t4}). Across the \TFourPairedN{}
settings where both methods return a feasible molecule, COMPOSE improves the
paired mean docking score by $0.335$ kcal/mol (95\% CI $0.019$--$0.651$ in favor
of COMPOSE). This performance is obtained with 500 docking-oracle evaluations per
setting, compared with 3{,}000 evaluations underlying the published GenMol
result, a sixfold reduction. Since many individual docking differences are
comparable to replicate variation, we use the paired aggregate as the primary
comparison rather than interpreting small cell-level differences separately.

\subsection{Capabilities beyond fixed-objective endpoint optimization}
\label{sec:gem-exp-capabilities}

Endpoint benchmarks score the molecule returned after optimization, but do not
capture how decisions are made or revised along the trajectory. Because every
COMPOSE state is a complete molecule, control can use downstream reachability,
resume from a realized intermediate after an objective change, or restrict which
molecular states may be visited.

\paragraph{Future reachability changes multi-step edit decisions.}
An edit that improves immediate similarity can still eliminate every successful
continuation within the remaining budget. We test this on 65 held-out
source--target pairs, each with a certified executable path to the target and
sources disjoint from development. The local and future-aware policies share the
same frozen $R_\theta$, target, edit budget, and executable successors. They
differ only in whether the next edit is selected by immediate target similarity
or by the continuations available over the remaining budget.
Local selection reaches 26 of 65 targets, whereas future-aware control reaches
40, rescuing 14 of the 39 local failures without losing a local success
(\cref{tab:mechanism}). In one deterministic rescue, the locally preferred edit
has higher immediate similarity but removes the successful continuation, while a
slightly less similar edit preserves a path to the target. Exhaustive lookahead
is not required: using $h_\phi$ to prioritize a single successor preserves the
40/65 recovery rate while evaluating 34.6\% of the continuations required by
all-candidate lookahead. The full prioritization comparison is reported in
\cref{app:exp-lookahead}.

\paragraph{Realized molecular states can be retargeted after an objective change.}
Starting from $x_0$, we execute three edits under an initial objective and then
switch objectives at the realized state $x_3$. We compare continuing from $x_3$
with restarting from $x_0$, holding the new objective and remaining edit budget
fixed.
Across 40 held-out molecules, continuing from $x_3$ improves post-switch utility
in both directions: $0.302\,[0.186,0.413]$ for the potency-first objective and
$0.176\,[0.069,0.283]$ for the developability-first objective
(\cref{tab:mechanism}). The continued trajectories also move toward the new
objective rather than preserving the preference under which $x_3$ was reached.
Retargeting requires no update to $R_\theta$; control resumes from the realized
molecular state under the new objective.

\paragraph{Pathwise constraints act on the molecular trajectory.}
Some molecular requirements must hold at every intermediate rather than only for
the returned endpoint. COMPOSE enforces such statewise constraints by removing
violating successors before the reference kernel is normalized. We compare this
support restriction with endpoint-only filtering under matched sources,
objectives, and edit budgets.
On the developmental panel, 16 of 24 endpoint-feasible trajectories contained at
least one intermediate constraint violation, showing that endpoint feasibility
can hide inadmissible paths. Support-restricted control excludes such
intermediates by construction. We report the matched effect on optimization
performance and the full protocol in \cref{app:exp-pathwise}.


